\documentclass[11pt]{article} 
\usepackage{nopageno}
\setcounter{secnumdepth}{1} % Use 1 for sections, 2 for subsections, etc.
\usepackage[margin=1in]{geometry}
\begin{document}

% Remove page numbers---Research.gov will paginate the document instead
%
\pagestyle{empty} 
\vspace{-0.3cm}
\subsection{Overview} 
\vspace{-0.3cm}
This project addresses the challenges faced by flood-vulnerable, heritage-rich inland non-metropolitan community centers (INMCCs) in preserving cultural heritage while recovering from increasingly frequent and severe flood events.
Building on our preliminary work, we introduce and develop the concept of the ``preservation-urgency paradox," which encapsulates the tension between rapid community mobilization and heritage preservation goals in post-flood recovery processes. 
Focusing on two case studies, Montpelier, VT and Marshall, NC, the project examines this paradox through a mixed-methods approach. 
The study aims to collect data on community capitals and engineering solutions for heritage recovery, analyze the interplay between community capitals and preservation efforts, and develop a ``Flood-Resilient Heritage Preservation Toolkit" for practical application in vulnerable communities.
\vspace{-0.5cm}
\subsection{Intellectual Merit} 
\vspace{-0.3cm}
Central to this study is the introduction and investigation of the ``preservation-urgency paradox," a concept developed from our preliminary work that challenges conventional assumptions about the relationship between community capitals and heritage preservation during inland flooding recovery. 
By examining this newly identified paradox, the study contributes to fundamental social science by advancing our understanding of community resilience, community capitals, and institutional adaptation in the context of heritage preservation and disaster recovery. 
The comparative analysis of multiple INMCCs at different stages of flood recovery will provide insights into how the preservation-urgency paradox manifests across diverse contexts, illuminating the ways that institutional arrangements and community capitals shape recovery outcomes. 
This work expands on community capital theory by exploring its nuanced role in disaster contexts, particularly in small, heritage-rich communities. 
Additionally, the research advances our understanding of knowledge integration processes, examining how different forms of expertise--local, technical, and institutional--are mobilized and synthesized in post-disaster decision-making. 
The development of the ``Flood-Resilient Heritage Preservation Toolkit" represents both an academic and practical contribution, offering a data-driven approach to addressing the complex challenges at the intersection of heritage preservation, disaster recovery, and climate adaptation, with specific strategies for navigating the preservation-urgency paradox in vulnerable communities.
\vspace{-0.5cm}
\subsection{Broader Impacts} 
\vspace{-0.3cm}
By addressing the challenges faced by flood-vulnerable, heritage-rich INMCCs through the lens of the newly developed ``preservation-urgency paradox,"  this project contributes to more equitable, effective, and culturally sensitive approaches to flood recovery and heritage preservation. 
The investigation of this paradox will yield insights that help preserve cultural identity, support local economies, and enhance community resilience in often-overlooked INMCCs facing increasing climate threats. 
The project maximizes its broader impacts through stakeholder engagement workshops that explore the manifestation of the preservation-urgency paradox in real-world settings and wide dissemination of the ``Flood-Resilient Heritage Preservation Toolkit" via collaborations with professional organizations. 
Training and education through webinars and conference presentations will focus on strategies for navigating the tension between urgent recovery needs and long-term preservation goals.
The real-world application of the toolkit through the Association for Preservation Technology's Disaster Response Initiative post-disaster field deployment program ensures immediate practical impact in addressing the preservation-urgency paradox. 
The research will inform policy and practice for more resilient and equitable recovery outcomes in flood-vulnerable communities, providing decision-makers with tools to balance rapid recovery imperatives with heritage preservation objectives. 
By integrating this work into engineering education, the project will foster interdisciplinary skills in future engineers, promoting a more holistic approach to disaster resilience and heritage preservation. 
\end{document}